
\section{Detailed Collected Data Sources}
\label{app:collect data sources}


\begin{table*}[htbp]
\centering
\footnotesize
\caption{Collected Telemetry by OS and Execution Environment}
\label{tab:collected_data_os}
\small
\begin{tabular}{@{}p{2.5cm} p{4cm} p{4cm} p{4cm}@{}}
\toprule
\textbf{Environment} &
\textbf{Process Telemetry} &
\textbf{Network Telemetry} &
\textbf{System/Auth Telemetry} \\
\midrule

\textbf{Windows} & 
Sysmon (process, registry, file events); PowerShell logs; WMI Activity; Task Scheduler; AppLocker; Application Experience &
Sysmon network events; Windows Firewall logs; DNS Client logs &
Windows Event Logs; RDP Core logs; Security Audit Configuration; AMSI \footnotemark\\
\midrule

\textbf{Linux} &
Sysmon for Linux (process, file, network events); Auditd &
Zeek telemetry (connections, DNS, HTTP, files); Suricata IDS telemetry (alerts, flows, DNS, HTTP, TLS) &
Syslog; PAM authentication logs; Audit logs \\
\midrule

\textbf{Containers} &
Tracee runtime telemetry (syscalls, probes, capability use, process events); Container runtime events &
Tracee network events &
cgroup metadata; mountinfo; docker API events \\
\bottomrule
\end{tabular}
\end{table*}

\footnotetext{
See Table~\ref{tab:windows_extended} for Windows Custom Channels.
}


\begin{table}[t]
\footnotesize
\setlength{\tabcolsep}{2pt}
\centering
\caption{Extended Windows Operational Channels and Associated Threat Indicators}
\label{tab:windows_extended}
\begin{tabular}{@{}lp{3.6cm}@{}}
\toprule
\textbf{Channel} & \textbf{Threat Indicators} \\
\midrule

PowerShell / Operational 
& Obfuscation, encoded payloads, script abuse \\

WMI-Activity / Operational 
& Fileless exec, lateral movement, WMI persistence \\

Security-Audit-Config-Client 
& Brute force, privilege escalation \\

AppLocker (MSI/Script; EXE/DLL) 
& Block/allow malicious binaries or scripts \\

TaskScheduler / Operational 
& Scheduled-task persistence \\

RdpCoreTS / Operational 
& RDP brute force, lateral movement \\

Program-Telemetry 
& Unknown or suspicious binaries \\

DNS-Client / Operational 
& Beaconing, C2 domains, DNS exfiltration \\

\bottomrule
\end{tabular}
\end{table}


\paragraph{Telemetry Availability and Variability}

We note that the availability and granularity of collected telemetry may vary across scenarios and execution environments. This variability is not an artifact of incomplete instrumentation, but a consequence of differences in attack configurations, execution outcomes, and platform-specific constraints. 

In particular, failed or partially executed attacks may only manifest as control-plane or network-level signals (e.g., authentication failures or connection attempts), without generating observable process creation, file system, or privilege-related events on the target host. Similarly, fileless execution techniques, short-lived processes, or container-scoped attacks may evade certain host-based sensors while remaining visible through alternative telemetry sources.

This design reflects real-world operational conditions, where defenders rarely observe a complete and uniform set of signals for every adversarial action. Rather than enforcing artificial completeness, our dataset preserves these natural observability gaps, enabling systematic analysis of when and why single-source telemetry fails, and how multi-source telemetry can mitigate such limitations. 

In scenarios where target-side Windows telemetry is sparse (e.g., due to partial execution, short-lived activity, or channel unavailability), we additionally extract attacker-side logs (e.g., Linux syslog) from the orchestration host. This attacker-side telemetry is used to provide auxiliary context for action provenance (e.g., command issuance, tooling invocation, and network attempt timing), rather than to replace host evidence on the victim. We explicitly preserve the trust-domain boundary between attacker-side and victim-side observations, and annotate each record with its collection origin (attacker vs.\ target) to support controlled analyses and to prevent inadvertent leakage of privileged signals into detection-only evaluations.

\paragraph{Host and Container Telemetry}

Windows systems contribute heterogeneous host-level telemetry from multiple independent Windows instrumentation sources, including Sysmon~\footnote{https://learn.microsoft.com/en-us/sysinternals/downloads/sysmon}, Windows Security Auditing, WMI Activity logs, and PowerShell operational logs. There sources are generated by distinct subsystems and logging pipelines, providing complementary visibility into process execution, authentication, management-plane activity, and script-level behavior. In addition to baseline system and security events, we collect extend extended operational channels (Table~\ref{tab:windows_extended}), which provide high-value visibility into adversarial behaviors such as command obfuscation, scheduled-task persistence, lateral movement, and DNS-based command-and-control activity.

Linux hosts provide complementary signals through Syslog, authentication logs (e.g., PAM), and network-layer metadata recorded by Zeek\footnote{https://zeek.org/} and Suricata\footnote{https://suricata.io/}. These network sensors capture DNS queries, HTTP transactions, protocol fingerprints, and intrusion alerts.

Containerised environments contribute additional system-level metadata, including mount-namespace structures, cgroup hierarchies, Docker socket exposure, and overlay filesystem metadata, offering visibility into execution patterns characteristic of CI/CD and supply-chain attack surfaces.

\paragraph{Behavioral Tracing and Enrichment}

To capture fine-grained behavioral signals beyond baseline system logs, we incorporate multiple system-call–level and kernel-level tracing mechanisms.

On Linux, Auditd and Sysmon-for-Linux provide detailed records of process creation, file interactions, network connections, and system-call sequences, allowing reconstruction of high-resolution behavioral traces. We also collected performance counters, but found them insufficient for reasoning about attack causality.

For containerised workloads, we additionally employ Tracee\footnote{https://www.aquasec.com/products/tracee/}, an eBPF-based runtime tracing framework. Tracee captures ephemeral kernel events—including privilege-escalation attempts, anomalous syscalls, and memory-resident execution—without requiring in-container agents. Its low-overhead, high-fidelity instrumentation is particularly effective for surfacing stealthy or short-lived behaviors common in supply-chain exploitation.

Together, these tracing mechanisms enrich the dataset with consistent, fine-grained behavioral signals across hosts and containers. These data are either ingested into Log Analytics Workspace directly or collected directed from hosts.
